\section*{ОБОЗНАЧЕНИЯ И СОКРАЩЕНИЯ}

БД -- база данных.

ИС -- информационная система.

ПО -- программное обеспечение.

СУБД -- система управления базами данных.

ТЗ -- техническое задание.

ТП -- технический проект.

UML (Unified Modeling Language) -- язык графического описания для объектного моделирования в области разработки программного обеспечения.

ER-модель -- модель данных, отражающая сущности предметной области и связи между ними.

HTTP (HyperText Transfer Protocol) -- протокол передачи данных между браузером пользователя и web-сервером.

HTML (HyperText Markup Language) -- язык гипертекстовой разметки web-страниц.

CSS (Cascading Style Sheets) -- каскадные таблицы стилей, используемые для оформления web-интерфейса.

PHP -- серверный язык программирования, используемый для обработки запросов web-приложения.

SQL (Structured Query Language) -- язык структурированных запросов для работы с базами данных.

PDO (PHP Data Objects) -- механизм PHP для подключения к базе данных и выполнения SQL-запросов.